\documentclass[../main.tex]{subfiles}

\begin{document}
    \section{SYSTEM DESIGN AND ARCHITECTURE}
    
    \subsection{Block Diagram}

    \vspace{2em}
    
    \begin{figure}[h]
        \centering
        \includegraphics[width=0.8\textwidth]{Images/block diagram.drawio.png}
    \vspace{0.2cm}
        \caption{Block diagram of our system}
        \label{fig:Block diagram}
    \end{figure}

   \subsubsection{Sensing Layer (Data Acquisition)}
   At the first level, multiple sensors are deployed within the poultry environment to continuously monitor critical parameters:
   \begin{itemize}
    \item DHT Sensor: Measures temperature and humidity to maintain optimal living conditions for poultry.
    \item Ammonia Sensor (MQ-137): Detects ammonia levels, which are harmful if they exceed safe limits.
    \item CO₂ Sensor (MH-Z19): Monitors carbon dioxide concentration to ensure proper ventilation.
    \item Lux Sensor: Measures light intensity, which affects poultry growth and behavior.
    \item Weight Sensor (Load Cell): Tracks the weight of poultry to monitor growth and productivity.
\end{itemize}

\subsubsection{ Processing Layer(Control Unit)}
All sensor data is transmitted to the  Arduino Nano, which acts as the central processing unit of the system.
\begin{itemize}
    \item The Arduino Nano collects and processes incoming sensor data.
    \item It performs basic filtering and validation of the data.
    \item Based on predefined threshold values, it determines whether the environmental conditions are within safe limits.
\end{itemize}
If abnormal conditions are detected, the Nano is designed to trigger appropriate actions such as:
\begin{itemize}
    \item Activating ventilation systems (fans)
    \item Controlling heating elements
    \item Generating alert signals
\end{itemize}

\subsubsection{ Communication Module}
An integrated communication module (Wi-Fi/LTE) enables the Nano to transmit processed data to a remote cloud server.
\begin{itemize}
    \item This ensures efficient and real-time data transfer with minimal latency.
\end{itemize}

\subsubsection{Cloud and Data Storage Layer}
The transmitted data is stored in a cloud database, where it is:
\begin{itemize}
    \item Logged for historical analysis
    \item Processed for visualization
    \item Used for generating alerts and insights
\end{itemize}

\subsubsection{Application Layer (User Interface)}
A web/mobile application dashboard provides an interface for the user (farmer) to interact with the system.Through this dashboard, the user can:
\begin{itemize}
    \item View real-time environmental data
    \item Analyze historical trends
    \item Receive alerts when parameters exceed safe thresholds
\end{itemize}


\subsubsection{Power Supply System}
The entire system is powered by a low power household grid supply along with protection circuit for voltage fluctuations.
\begin{itemize}
    \item Improves reliability and sustainability of the system
\end{itemize}

\newpage
\subsection{Flowchart}

 \begin{figure}[h]
        \centering
        % \includegraphics[scale=0.41]
        \includegraphics[width=0.6\textwidth]{Images/poultryflowchart.drawio (1).png}
        \caption{Flowchart for overall system} 
        \label{fig:overall_system}
    \end{figure}




    \subsubsection{System Initialization}
    The process begins with the start of the system, where all hardware components are powered on.
    \begin{itemize}
        \item All sensors are initialized and calibrated to ensure accurate readings.
        \item The Arudino Nano establishes connection with the network (Wi-Fi/LTE).
    \end{itemize}

    \subsubsection{Data Collection}
    After initialization:
   \begin{itemize}
       \item The system continuously collects data from all connected sensors.
       \item Parameters include temperature, humidity, ammonia, CO₂, light intensity, and weight.
   \end{itemize}

   \subsubsection{Data Transmission and Processing}
   The collected data is transmitted to the cloud/server using IoT protocols and the system validates the data by checking for:
   \begin{itemize}
       \item Errors
        \item Missing values
        \item Redundant data
   \end{itemize}

\subsubsection{Decision Making}
The system evaluates:
\begin{itemize}
    \item Threshold values are predefined for each parameter.
    \item This is the core decision point of the system.
\end{itemize}

\subsubsection{If Conditions are NORMAL}
If all parameters are within safe limits:
\begin{itemize}
    \item Data is stored in the database
    \item Data is displayed on the web/mobile dashboard
\end{itemize}

\subsubsection{If Conditions are ABNORMAL}
If any parameter exceeds the threshold:
\begin{itemize}
    \item The system identifies the specific parameter causing the issue
    \item It generates alerts (notifications to the farmer)
    \item It may also trigger corrective actions, such as:
\begin{itemize}
    \item Turning on fans
     \item Activating heaters
      \item Adjusting ventilation
\end{itemize}
\end{itemize}

\subsubsection{Continuous Monitoring Loop}
\begin{itemize}
    \item The entire process runs in a continuous loop
    \item Ensures real-time monitoring and rapid response
\end{itemize}

    
\newpage

\end{document}
